%%%%%%%%%%%%%%%%%%%%%%%%%%%%%%%%%%%%%%%%%%%%%%%%%%%%%%%%%%%%%%%%%%%%%%%%
% Plantilla TFG/TFM
% Escuela Politécnica Superior de la Universidad de Alicante
% Realizado por: Jose Manuel Requena Plens
% Contacto: info@jmrplens.com / Telegram:@jmrplens
%%%%%%%%%%%%%%%%%%%%%%%%%%%%%%%%%%%%%%%%%%%%%%%%%%%%%%%%%%%%%%%%%%%%%%%%

\chapter{Anexo - Código fuente del sistema}

El código fuente completo del sistema desarrollado en este Trabajo Final de Máster se encuentra disponible en un repositorio público alojado en la plataforma GitHub. En dicho repositorio se incluye tanto el backend del sistema, encargado de la comunicación con los dispositivos de red, como el frontend web utilizado para la visualización y gestión de la información de monitorización.

El repositorio contiene la estructura completa del proyecto, así como los ficheros necesarios para su despliegue y ejecución en un entorno local. El acceso al código se proporciona con fines de consulta y revisión, sin que sea necesario su análisis detallado para la comprensión del presente documento.

El repositorio puede consultarse en la siguiente dirección:

\begin{center}
\url{https://github.com/metalolmo/TFM-monitorizacion-red}
\end{center}

En el momento de la entrega de este trabajo, la rama principal del repositorio corresponde a la versión estable del sistema descrito en la memoria.